\documentclass[12p]{article}

\usepackage{bbm}
\usepackage{fancyhdr}
\usepackage{titling}
\usepackage{enumitem}
\usepackage{amsmath}
\usepackage{dsfont}
\usepackage{amsfonts}
\usepackage{amssymb}
\usepackage{amsthm}
\usepackage{geometry}
\usepackage{tikz}
\usepackage{algorithm}
\usepackage{algpseudocode}
\geometry{verbose,tmargin=1in,bmargin=1in,lmargin=1in,rmargin=1in}
\renewcommand{\vec}[1]{\mathbf{#1}}
\newcommand{\dec}{\text{dec}}
\newcommand{\divi}{\mathbin{\textnormal{div}}}
\title{Homework 1}
\author{Satvik Prasad}
\date{}
\newtheorem{lemma}{Lemma}[section]
\fancypagestyle{plain}{%
  \fancyhf{}
  \fancyhead[L]{\theauthor}
  \fancyhead[R]{\thedate}
}
\pagestyle{plain}
\makeatletter
\def\@maketitle{%
  \newpage
  \null
  \vskip 1em%
  \begin{center}%
    \let \footnote \thanks
    {\LARGE \@title \par}%
    \vskip 1em%
  \end{center}%
  \par
\vskip 1em}
\makeatother
\usepackage{cmbright}

\newcommand{\mat}[1]{\mathbf{#1}}

% ---------- Probability theory ----------
% Core operators (optional argument-free; use as \E{X}, \Var{X}, \Cov{X}{Y})
\DeclareMathOperator{\EOp}{\mathbb{E}}
\DeclareMathOperator{\VarOp}{Var}
\DeclareMathOperator{\CovOp}{Cov}
\DeclareMathOperator{\CorrOp}{Corr}

\renewcommand{\P}[1]{\mathbb{P}\left(#1\right)}        % \P{A}
\newcommand{\E}[1]{\EOp\left[#1\right]}                % \E{X}
\newcommand{\Var}[1]{\VarOp\left(#1\right)}            % \Var{X}
\newcommand{\Cov}[2]{\CovOp\left(#1,\,#2\right)}       % \Cov{X}{Y}
\newcommand{\Corr}[2]{\CorrOp\left(#1,\,#2\right)}     % \Corr{X}{Y}
\newcommand{\SD}[1]{\operatorname{SD}\left(#1\right)}  % \SD{X}

% Subscripted versions, e.g. \Esub{\theta}{X}, \Psub{\theta}{A}
\newcommand{\Esub}[2]{\EOp_{#1}\left[#2\right]}
\newcommand{\Psub}[2]{\mathbb{P}_{#1}\left(#2\right)}

% Conditioning, indicators, independence
\newcommand{\given}{\,\middle|\,}                      % \E{X \given Y}
\newcommand{\ind}[1]{\mathbbm{1}\left\{#1\right\}}     % \ind{X > 0}
\newcommand{\indep}{\perp\!\!\!\perp}                  % X \indep Y
\newcommand{\iid}{\overset{\text{iid}}{\sim}}          % X_i \iid \Normal(0,1)

% Convergence
\newcommand{\convp}{\xrightarrow{\;p\;}}               % in probability
\newcommand{\convd}{\xrightarrow{\;d\;}}               % in distribution
\newcommand{\convas}{\xrightarrow{\text{a.s.}}}        % almost surely
\newcommand{\convLp}[1]{\xrightarrow{L^{#1}}}          % in L^p

% Probability spaces
\newcommand{\F}{\mathcal{F}}                           % sigma-algebra
\newcommand{\B}{\mathcal{B}}                           % Borel sigma-algebra
\newcommand{\pspace}{(\Omega, \F, \mathbb{P})}

% Common distributions
\newcommand{\Normal}{\mathcal{N}}
\newcommand{\Bern}{\operatorname{Bernoulli}}
\newcommand{\Bin}{\operatorname{Binomial}}
\newcommand{\Pois}{\operatorname{Poisson}}
\newcommand{\Unif}{\operatorname{Uniform}}
\newcommand{\Exp}{\operatorname{Exp}}
\newcommand{\Geom}{\operatorname{Geometric}}
\newcommand{\Gam}{\operatorname{Gamma}}
\newcommand{\Beta}{\operatorname{Beta}}

% Number sets
\newcommand{\R}{\mathbb{R}}
\newcommand{\N}{\mathbb{N}}
\newcommand{\Z}{\mathbb{Z}}


\begin{document}
\begin{enumerate}
  \item $X \sim \text{Binomial}(N, p)$ is a sum of $N$ Bernoulli
    trials with probability $p$ of success, yielding generating function
    $$
    G_X(s) = (1-p + ps)^N.
    $$

    \begin{enumerate}
      \item Since $\mathbb E[X] = G_X'(1),$ we get
        $$
        G_X'(s) = Np(1-p + ps)^{N-1} \Rightarrow \mathbb E[X] = Np.
        $$
        Similarly, since $\mathbb E[X^2] - \mathbb E[X] = G_X''(1),$ we get
        $$
        N(N-1) p^2 = \mathbb E[X^2] - Np \Rightarrow \mathbb E[X^2] =
        N^2p^2 - Np^2 + Np,
        $$
        and $\text{Var}\,X = N^2p^2 - Np^2 + Np - N^2p^2 = Np(1-p)$ as required.

      \item The probability that $X$ is even can be computed from
        $$
        \frac{G_X(1) + G_X(-1)}{2},
        $$
        since odd exponents cancel out. Thus,
        $$
        \mathbb P(X \equiv 0 \pmod 2) = \dfrac{1^N + (1-2p)^N}{2} =
        \dfrac{1 + (1-2p)^N}{2}.
        $$

      \item Use roots of unity such that $\omega^2 + \omega + 1 = 0.$
    \end{enumerate}

  \item We have random variable $X$ with generating function
    $G_X(s).$ Consider $Y = kX$ for some positive integer $k.$ Then,
    $$
    G_{Y}(s) = \mathbb E[s^{kX}] = \mathbb E[(s^k)^X] = G_X(s^k).
    $$

    Similarly, if $Z = X + k,$ then $$
    G_Z(s) = \mathbb E[s^{X + k}] = s^k \mathbb E[s^X] = s^k G_X(s)
    $$
    as required. \qed

  \item The generating function for $u_n$ is given by
    \begin{align*}
      U(s) = \sum_{k = 0}^{\infty} \mathbb P(X > n) s^k &= \sum_{k =
      0}^{\infty} \sum_{j = k + 1}^{\infty} \mathbb P(X = j) s^k \\
      &=\sum_{j = 1}^{\infty} \mathbb P(X = j) \sum_{k = 0}^{j - 1} s^k \\
      &= \sum_{j =1 }^\infty{} \mathbb P(X = j) \dfrac{s^{j} - 1}{s - 1} \\
      &= \dfrac{1}{s-1} \sum_{j = 1}^{\infty} \mathbb P(X =j) s^j -
      \dfrac{1}{s-1} \mathbb P(X > 0) \\
      &= \dfrac{1}{s-1} (G_X(s) - \mathbb P(X = 0)) - \dfrac{1}{s-1}
      \mathbb P(X > 0) \\
      &= \dfrac{1}{s-1} G_X(s) - \dfrac{1}{s-1},
    \end{align*}
    or equivalently, $$
    (1-s) U(s) = 1 - G_X(s).
    $$ \qed

  \item Let $X$ be the random variable that describes the total
    number of successes in $N$ turns. Hence, $X$ will be the sum of
    $N$ random variables describing the total number of successes in
    1 turn. Let this random variable be $Y.$ Then,
    $$
    G_Y(s) = \sum_{k = 0}^{\infty} p^{k} (1-p) s^{k} = \dfrac{1-p}{1-ps}.
    $$
    Thus, $$G_X(s) = \dfrac{(1-p)^N}{(1-sp)^N}.$$

    To find the mean, we get
    \begin{align*}
      G'_X(s) = \dfrac{pN(1-p)^{N}}{(1-sp)^{N + 1}} \Rightarrow
      \mathbb E[X] = \dfrac{Np}{1-p}
    \end{align*}
    as required.

    Differentiating a second time, we get
    $$
    G_X''(s) = \dfrac{p^2N(N+1)(1-p)^{N}}{(1-sp)^{N+2}} \Rightarrow
    \mathbb E[X^2] - \mathbb E[X] =  p^2 N(N+1) (1-p)^{-2},
    $$
    meaning
    $$
    \mathbb E[X^2] = p^2N(N+1) (1-p)^{-2} + Np(1-p)^{-1},
    $$
    i.e.,
    \begin{align*}
      \text{Var}\, X &= p^2N(N+1)(1-p)^{-2} + Np(1-p)^{-1} - N^2p^2(1-p)^{-2} \\
      &= p^2N(1-p)^{-2} + Np(1-p)^{-1} \\
      &= \dfrac{Np}{(1-p)^2}
    \end{align*}
    as required. \qed

  \item
    \begin{enumerate}
    \item From 1a), the answer is $\mathbb E[X] = np$ and
      $\text{Var}\, X = np(1-p).$
  \item From 1b), the answer is
    $$
    \dfrac{1 + (1-2p)^{n}}{2}.
    $$
  \item Consider
    $$
    \dfrac{G_X(1) + G_X(\omega) + G_X(\omega^2)}{3}
    $$
    where $\omega^3 = 1, \omega \in \mathbb C \setminus \{1\}.$ Since
    $\omega^2 + \omega + 1 = 0,$ only terms where $k \equiv 0 \pmod
    3$ survive, meaning the above is $\mathbb P(3 \mid X).$

    Substituting, we get
    $$
    \mathbb P(3 \mid X) = \dfrac{1 + (1 -p + p\omega)^{n} + (1-p +
    p\omega^2)^n}{3}.
    $$
    Since
    $$
    \omega^2 = \bar \omega,
    $$
    we get
    \begin{align*}
      (1 - p + p\omega)^n + (1- p + p \omega^2)^n &= (1- p +
      p\omega)^n + (1 - p + p \bar\omega)^n \\
      &=2 \Re (q  + p\omega)^n.
    \end{align*}
    Taking $\omega = (-1 + \sqrt3 i)/2,$ we get
    $$
    1 - p + p\omega = 1 - \dfrac{3p}{2} + \dfrac{p\sqrt{3}}{2} i.
    $$
    Hence,
    $$
    |1 - p + p \omega| = \sqrt{\left(1 - \dfrac{3p}{2}\right)^2 +
    \dfrac{3p^2}{4}} = \sqrt{1 - 3p + 3p^2}.
    $$
    Let
    $$
    \theta = \arg\left(1 - \dfrac{3p}{2} + \dfrac{p\sqrt3}{2} i\right),
    $$
    so that $1 - p + p\omega = \sqrt{1 - 3p + 3p^2}\, e^{i\theta}$
    and, by De Moivre's theorem,
    $$
    (1 - p + p\omega)^n = \left(1 - 3p + 3p^2\right)^{n/2} e^{in\theta}.
    $$
    Taking real parts,
    $$
    \Re (1 - p + p\omega)^n = \left(1 - 3p + 3p^2\right)^{n/2} \cos(n\theta).
    $$
    Therefore,
    $$
    \mathbb P(3 \mid X) = \dfrac{1 + 2\left(1 - 3p +
    3p^2\right)^{n/2} \cos(n\theta)}{3}.
    $$
    \qed
\end{enumerate}
\end{enumerate}
\end{document}
